% ===========================================================================
% COL334 Assignment 2 -- The Socket Exchange
% Experiment Report (handout section 8) + Bonus (handout section 6.9)
%
% Build:  pdflatex report.tex   (run twice for the table of contents)
% ===========================================================================
\documentclass[11pt,a4paper]{article}

\usepackage[margin=2.2cm]{geometry}
\usepackage{graphicx}
\usepackage{booktabs}
\usepackage{array}
\usepackage{fancyvrb}
\usepackage[hidelinks]{hyperref}
\usepackage{caption}

\captionsetup{font=small,labelfont=bf}
\setlength{\parskip}{0.5em}
\setlength{\parindent}{0pt}

\newcommand{\shot}[2]{%
  \begin{figure}[htbp]
    \centering
    \fbox{\includegraphics[width=\textwidth]{#1}}
    \caption{#2}
  \end{figure}%
}

\newcommand{\cmd}[1]{\texttt{\small #1}}

\title{\textbf{The Socket Exchange}\\[0.3em]
       \large COL334 --- Assignment 2 --- Experiment Report}
\author{Rohit Meena (2024CS10030) \quad\textbullet\quad Subarno Saha (2024CS50431)}
\date{FreeBSD 14.4-RELEASE \quad\textbullet\quad C++17 \quad\textbullet\quad \today}

\begin{document}
\maketitle

All experiments were run on FreeBSD 14.4-RELEASE. The Exchange Server is written
in C++17 against the POSIX socket API with \cmd{kqueue()}; no third-party
networking library is used. Every measurement below was taken on the running
system with the FreeBSD tools named in each section.

\tableofcontents
\newpage

% ===========================================================================
\section{Implementation Decisions}
% ===========================================================================

\subsection{Concurrency / I/O approach}

The Exchange Server is a \textbf{single-threaded event loop built on
\cmd{kqueue()}}. There is one process and one thread for the entire server. The
listening socket and every accepted connection are registered on a single kqueue
with \cmd{EVFILT\_READ}; \cmd{kevent()} blocks until the kernel reports
descriptors ready, and the loop services only those.

This is verifiable from outside the process. \cmd{procstat -f} shows the server
holding the listening socket on fd~3, \textbf{the kqueue itself on fd~4}
(descriptor type \cmd{k}), and one socket fd per connection thereafter
(Figures~\ref{fig:e1b} and~\ref{fig:e5b}). \cmd{procstat -k} shows exactly one
thread, asleep inside \cmd{sys\_kevent}:

\begin{Verbatim}[fontsize=\small,frame=single]
mi_switch sleepq_catch_signals sleepq_wait_sig _sleep kqueue_scan
kqueue_kevent kern_kevent_fp kern_kevent_generic sys_kevent
amd64_syscall fast_syscall_common
\end{Verbatim}

\subsection{Why this approach}

\textbf{Cost scales with activity, not with connection count.} \cmd{kevent()}
returns only descriptors the kernel has already marked ready, so a loop
iteration costs $O(\mathrm{ready})$ rather than $O(\mathrm{open})$. Experiment~5
shows this directly: with five connections open and only three carrying data,
the server called \cmd{recv()} on exactly those three and never touched the
other two.

\textbf{One unresponsive client cannot stall the others.} Because the loop
blocks in \cmd{kevent()} rather than in \cmd{recv()} on a particular socket, a
client that connects and then sends a partial message is simply never reported
ready. Experiment~4 measured a second client being served in
0.001\,s while the first held an unterminated message, with
\cmd{procstat -k} confirming the server was parked in \cmd{sys\_kevent} and not
in \cmd{soreceive}.

\textbf{Per-connection cost is tiny, which is what makes the bonus possible.}
Measured across seven levels, server resident memory grows by
$\approx$\,2012\,KB per 10{,}000 idle connections --- about
\textbf{201 bytes per connection} --- on a 3.4\,MB baseline. Holding
70{,}000 idle connections cost 17.1\,MB of RSS and 0.1\% CPU.

\textbf{The alternatives are worse.} A thread per connection would need 70{,}000
stacks; even at a conservative 512\,KB each that is
$\approx$\,34\,GB against 4\,GB of RAM, roughly four
orders of magnitude more than the measured 201~bytes. \cmd{select()} is bounded
by \cmd{FD\_SETSIZE} (1024 by default) and so cannot reach these numbers at all.
\cmd{poll()} has no such cap but requires passing and scanning the whole
descriptor array on every call --- $O(n)$ in the number of \emph{open}
descriptors, i.e.\ 70{,}000 entries examined to find the few that are ready.

\subsection{Other significant design decisions}

\textbf{All sockets are non-blocking.} \cmd{set\_nonblocking()} is applied to
every descriptor returned by \cmd{accept()}. This is mandatory rather than
optional: in a readiness-based loop a single blocking \cmd{recv()} or
\cmd{send()} would stall every other client, which is exactly the failure
Experiment~4 tests for.

\textbf{Message framing is explicit, because TCP provides none.} Each connection
owns a \cmd{LineBuffer} that accumulates bytes from \cmd{recv()} and yields
complete \verb|\n|-terminated lines. One \cmd{recv()} is never assumed to equal
one message. Experiment~3 demonstrates why: a single 24-byte message sent in
four writes arrived as four separate \cmd{recv()} calls of 6, 10, 7 and 1 bytes,
and the server produced no command until the final newline arrived. On the
output side the \verb|\n| is appended in exactly one place (\cmd{enqueue()}), so
a line can never be double-framed.

\textbf{Backpressure is handled per connection, so a slow reader is contained.}
\cmd{enqueue()} first attempts \cmd{send()} directly. If the peer's receive
window is full, \cmd{send()} returns \cmd{EAGAIN} --- treated as backpressure,
not as an error --- and the unsent remainder is appended to that connection's
own \cmd{outbuf} queue, with \cmd{EVFILT\_WRITE} registered for that socket
only. When the socket becomes writable the queue is drained and the filter is
removed again. Consequently a client that stops reading backs up on its own
socket alone. Experiment~7 confirms this: with one market-data client
deliberately not reading, its receive queue grew to 85{,}000 bytes while the
normal client stayed at 17~bytes and both traders completed all 5{,}000 trades.

\textbf{The accept loop drains fully.} A single \cmd{EVFILT\_READ} on the
listening socket can represent several queued connections, so \cmd{accept()} is
called in a loop until it returns \cmd{EAGAIN}. Without this the server would
fall behind under burst load --- which matters at the connection rates used in
the bonus.

\textbf{Connection teardown is idempotent.} \cmd{close\_session()} returns
immediately if the session is already gone, and \cmd{EV\_EOF} is acted on only
after \cmd{on\_readable()} has run, so a descriptor is never closed twice.
\cmd{recv()} returning 0 (FIN) and \cmd{recv()} failing with \cmd{ECONNRESET}
(RST) are logged distinctly, which is what made Experiments~6 and~8 measurable
from the application side.

\textbf{Roles are inferred from the first command.} A connection is
\cmd{UNKNOWN} until it issues a command; \cmd{LOGIN}/\cmd{BUY}/\cmd{SELL}/%
\cmd{CANCEL} make it a Trader, \cmd{SUBSCRIBE}/\cmd{UNSUBSCRIBE} a Market-Data
client. A command outside the connection's role is answered with \cmd{ERROR},
which is what makes a Market-Data connection genuinely read-only.

\textbf{Resting orders outlive their owner's connection} (handout 2.6).
\cmd{close\_session()} deliberately does \emph{not} remove the departing
client's orders; they remain in the book and can still match. Subscribed
Market-Data clients still receive the resulting \cmd{TRADE} messages, while the
departed trader receives no \cmd{BOUGHT}/\cmd{SOLD}.

\textbf{A generation counter guards against descriptor reuse.} Every session
carries a monotonically increasing \cmd{session\_seq}, recorded on each order.
Before delivering a fill the server checks that the descriptor's \emph{current}
session sequence still matches the one that submitted the order. Without this, a
fill belonging to a departed trader could be delivered to whichever new client
happened to inherit that file descriptor number --- a real hazard given how
aggressively descriptors are recycled.

\newpage
% ===========================================================================
\section{Experiments}
% ===========================================================================

% ---------------------------------------------------------------------------
\subsection{Experiment 1 --- Listening and Connected Sockets}

\textbf{Question.} After the client has connected, identify the TCP sockets
associated with the Exchange Server. How does the listening socket differ from
the socket representing the connection to the client?

\textbf{Answer.} The server holds \textbf{two distinct TCP sockets on two
distinct file descriptors}. Fd~3 is the \emph{listening} socket: bound to the
local half-association \cmd{127.0.0.1:5000} only, with a wildcard foreign
address \cmd{*:*}, in state \cmd{LISTEN}. It carries no peer and no data and
exists solely to be \cmd{accept()}ed from. Fd~5 is the \emph{connected} socket
returned by that \cmd{accept()}: fully specified by the four-tuple
\cmd{127.0.0.1:5000 $\leftrightarrow$ 127.0.0.1:51080}, in state
\cmd{ESTABLISHED}, and the only one of the two over which bytes flow. Both share
local port 5000; it is the \emph{remote} endpoint that distinguishes them, which
is precisely why one listening port can fan out to many simultaneous
connections.

\textbf{Approach.} Ran the harness so the server was listening with exactly one
client attached, then inspected the same instant from two angles: the
system-wide socket table (which addresses and states exist) and the server
process's own descriptor table (which fds those sockets occupy inside the
process).

\textbf{Commands and tools.}
\begin{Verbatim}[fontsize=\small,frame=single]
sockstat -4l | grep 5000                    # baseline: port free before start
sockstat -4  | grep -E 'USER|exchange|5000' # server sockets by process and fd
netstat -an -p tcp | grep -E 'Proto|5000'   # addresses and TCP states
procstat -f $(pgrep -x exchange_server)     # the server's fd table
\end{Verbatim}

\textbf{Evidence.} \cmd{sockstat} shows fd~3 $\rightarrow$ \cmd{*:*} and fd~5
$\rightarrow$ \cmd{127.0.0.1:51080}. \cmd{netstat} shows one \cmd{LISTEN} row
and \emph{two} \cmd{ESTABLISHED} rows (both directions of the loopback pair).
\cmd{procstat -f} additionally shows fd~4 of type \cmd{k} --- the kqueue.
Recv-Q and Send-Q are 0 on every row: an idle connection buffers nothing.

\shot{/home/vzux/Desktop/A2_evidence/exp1/exp1_listen-vs-established.png}
{Experiment 1 --- baseline showing the port free, then \texttt{sockstat} and
\texttt{netstat} with the \texttt{LISTEN} socket (\texttt{*:*}) and both
\texttt{ESTABLISHED} rows of the loopback pair.}

\shot{/home/vzux/Desktop/A2_evidence/exp1/exp1_server-fds.png}
{Experiment 1 --- \texttt{procstat -f}: fd~3 listening socket, \textbf{fd~4 the
kqueue} (type \texttt{k}), fd~5 the connected socket. One process holds all
three.}
\label{fig:e1b}

\newpage
% ---------------------------------------------------------------------------
\subsection{Experiment 2 --- Observing TCP Connection States}

\textbf{Question.} How does the TCP state of the client--server connection
change during its lifetime, and what events cause the observed state changes?

\textbf{Answer.}

\emph{Establishment.} \cmd{connect()} sends SYN; the server replies SYN+ACK; the
client ACKs. Both endpoints reach \cmd{ESTABLISHED}. On loopback this took
20\,$\mu$s (SYN at 06:14:20.569648, final ACK at .569668), so
\cmd{SYN\_SENT} and \cmd{SYN\_RCVD} never appear in a 1-second \cmd{netstat}
sample --- they exist only in the packet capture.

\emph{Active phase.} Both rows remain \cmd{ESTABLISHED} with
Recv-Q~=~Send-Q~=~0 for the full 10-second idle window.

\emph{Termination.} The client closes first, sending FIN at 06:14:30.658923. The
server ACKs it and its \cmd{recv()} returns 0 --- logged as
\cmd{[5] recv() -> 0 : peer closed (FIN)} --- whereupon the event loop closes
the socket, emitting the server's own FIN at .659202, which the client ACKs.

\emph{The critical timing result.} The gap between the two FINs is
\textbf{279\,$\mu$s} (Pass~A) and 143\,$\mu$s (Pass~B).
That interval \emph{is} the server's \cmd{CLOSE\_WAIT} and the client's
\cmd{FIN\_WAIT\_2} state. Both are real but last only microseconds, because the
event loop reacts to \cmd{recv()==0} immediately. Their absence from
\cmd{netstat} is a sampling artefact, not an absence of the states.

\emph{\cmd{TIME\_WAIT} and the loopback special case.} This guest has
\cmd{net.inet.tcp.nolocaltimewait = 1}, so FreeBSD deliberately skips
\cmd{TIME\_WAIT} when both endpoints are local. Hence in Pass~A the two rows
simply vanish at 06:14:31, leaving only \cmd{LISTEN}. Setting the sysctl to 0
(Pass~B) restores standard behaviour: the closing side enters \cmd{TIME\_WAIT}
at 06:36:41 and holds it through 06:37:40, disappearing at 06:37:41 ---
\textbf{exactly 60\,s $= 2 \times$ MSL}, with \cmd{net.inet.tcp.msl = 30000}\,ms.

\emph{\cmd{TIME\_WAIT} is kernel state, not process state.} The harness stopped
the server 15\,s after the close: the \cmd{LISTEN} row is present at 06:36:56 and
gone at 06:36:57, yet \cmd{TIME\_WAIT} survives a further \textbf{44 seconds}, to
06:37:40. The four-tuple is held by the kernel long after the owning process has
exited --- which is the purpose of \cmd{TIME\_WAIT}: preventing a delayed
duplicate segment from being delivered to a new connection reusing that
four-tuple. The contrast is direct: in Pass~A the \cmd{LISTEN} row dies
\emph{with} the process and nothing outlives it.

\emph{Incidental observation.} The harness's readiness probe reaches port 5000
before the server has bound, at 06:14:20.462364, and its SYN is answered
32\,$\mu$s later by RST --- connection refused. A SYN to a port with
no listener produces RST, never FIN. This is reused in Experiment~6.

\textbf{Approach.} Ran the experiment twice with three terminals running
concurrently: a packet capture started first, a 1-second \cmd{netstat} sampler
second, the harness last. Pass~A used the guest default; Pass~B set
\cmd{nolocaltimewait=0} to expose \cmd{TIME\_WAIT}, and the sampler was
deliberately left running past the end of the harness so the \cmd{TIME\_WAIT}
lifetime could be measured against the disappearance of the \cmd{LISTEN} row.

\textbf{Commands and tools.}
\begin{Verbatim}[fontsize=\small,frame=single]
su root -c 'tcpdump -l -i lo0 -n -tttt "tcp port 5000"' | tee exp2A_tcpdump.txt
while true; do date +%T; netstat -an -p tcp | grep 5000; echo; sleep 1; done \
    | tee exp2A_states.txt
python3 -u experiment.py 2 2>&1 | tee exp2A_harness.txt
su root -c 'sysctl net.inet.tcp.nolocaltimewait=0'   # Pass B; restored to 1 after
sysctl net.inet.tcp.msl                              # = 30000 ms
\end{Verbatim}
The \cmd{-l} flag is essential: without it \cmd{tcpdump} block-buffers into the
pipe and prints nothing until it exits.

\shot{/home/vzux/Desktop/A2_evidence/exp2/exp2_A-established.png}
{Experiment 2, Pass A --- steady \texttt{ESTABLISHED} (both directions) plus
\texttt{LISTEN}, 06:14:20--24.}

\shot{/home/vzux/Desktop/A2_evidence/exp2/exp2_A-after-close.png}
{Experiment 2, Pass A --- the transition. \texttt{ESTABLISHED} at 06:14:28 and
:30; at 06:14:31 both rows have vanished with \textbf{no} \texttt{TIME\_WAIT},
because \texttt{nolocaltimewait=1} skips it on loopback.}

\shot{/home/vzux/Desktop/A2_evidence/exp2/exp2_A-tcpdump-fin.png}
{Experiment 2, Pass A --- the complete 13-packet capture: the connection-refused
RST, the three-way handshake, and the four-way FIN close.}

\shot{/home/vzux/Desktop/A2_evidence/exp2/exp2_A-listen-dies-with-process.png}
{Experiment 2, Pass A --- \texttt{LISTEN} persists to 06:14:44 and disappears at
:45 when the server process exits. Nothing outlives the process.}

\shot{/home/vzux/Desktop/A2_evidence/exp2/exp2_B-sysctl-flipped.png}
{Experiment 2, Pass B --- \texttt{net.inet.tcp.nolocaltimewait: 1 -> 0}.}

\shot{/home/vzux/Desktop/A2_evidence/exp2/exp2_B-time-wait-outlives-process.png}
{Experiment 2, Pass B --- the key result. At 06:36:56 both \texttt{TIME\_WAIT}
and \texttt{LISTEN} are present; at 06:36:57 \texttt{LISTEN} is gone (the server
process has exited) yet \texttt{TIME\_WAIT} continues.}

\shot{/home/vzux/Desktop/A2_evidence/exp2/exp2_B-time-wait-expires-60s.png}
{Experiment 2, Pass B --- last \texttt{TIME\_WAIT} sample at 06:37:40, gone at
06:37:41: exactly 60\,s $=2\times$MSL after the close at 06:36:41.}

\shot{/home/vzux/Desktop/A2_evidence/exp2/exp2_B-harness-server-log.png}
{Experiment 2, Pass B --- server trace: accept, \texttt{recv() -> 0 : peer
closed (FIN)}, close.}

\newpage
% ---------------------------------------------------------------------------
\subsection{Experiment 3 --- TCP as a Byte Stream}

\textbf{Question.} Does the Exchange Server receive the application-level
message as one complete unit, or can it be received in multiple pieces? What
does this demonstrate about TCP and application-level message boundaries?

\textbf{Answer.} \textbf{In multiple pieces.} The single message
\cmd{LOGIN experiment\_trader\textbackslash n} (24 bytes) was written by the
client as four separate \cmd{send()} calls and arrived as \textbf{four separate
\cmd{recv()} calls returning 6, 10, 7 and 1 bytes} --- never as one 24-byte
unit. The server could act on none of them until the fourth delivered the
terminating newline, at which point it emitted
\cmd{[5] LOGIN experiment\_trader}.

This demonstrates that TCP is a \emph{byte stream}, not a message stream. It
guarantees that bytes arrive in order and without loss, but preserves no record
of where the sender's writes began or ended. Application-level message
boundaries do not exist at the transport layer, so the application must impose
them --- here by buffering and splitting on \verb|\n|. The converse also holds:
had the pieces arrived close enough together they would have coalesced into one
\cmd{recv()}, and several messages can equally arrive in one \cmd{recv()}.
Either way the receiver cannot assume one \cmd{recv()} equals one message, which
is why the \cmd{LineBuffer} accumulator is required rather than optional.

\textbf{Approach.} Observed the same event from both sides simultaneously: the
server's own instrumentation logging the return value of every \cmd{recv()}, and
a packet capture with ASCII payload decoding to confirm what actually crossed
the wire. Cross-checking the two rules out the possibility that the split was an
artefact of the logging rather than of TCP.

\textbf{Commands and tools.}
\begin{Verbatim}[fontsize=\small,frame=single]
su root -c 'tcpdump -l -i lo0 -n -tttt -A "tcp port 5000"' | tee exp3_tcpdump.txt
python3 -u experiment.py 3 2>&1 | tee exp3_harness.txt
\end{Verbatim}
\cmd{-A} prints payload as ASCII. The server logs each \cmd{recv()} return value
at \cmd{exchange\_server.cpp:365}; handout 6.3 explicitly permits modifying the
implementation to make this observable.

\textbf{Evidence.} Four \cmd{[P.]} segments from \cmd{127.0.0.1:28206}, with
\emph{contiguous} sequence numbers --- TCP delivered one unbroken 24-byte stream;
the segmentation reflects the sender's write pattern, not the message:

\begin{center}\small
\begin{tabular}{llrl}
\toprule
Time & Seq range & Length & Payload \\
\midrule
15:59:02.821997 & 1:7   & 6  & \cmd{"LOGIN "} \\
15:59:03.030565 & 7:17  & 10 & \cmd{"experiment"} \\
15:59:03.241641 & 17:24 & 7  & \cmd{"\_trader"} \\
15:59:03.451747 & 24:25 & 1  & \verb|"\n"| \\
15:59:03.452274 & \multicolumn{2}{l}{server $\rightarrow$ client, 3} & \cmd{"OK"} \\
\bottomrule
\end{tabular}
\end{center}

The server's \cmd{OK} left 527\,$\mu$s after the newline arrived,
confirming it was waiting for the delimiter and not for a message.

\shot{/home/vzux/Desktop/A2_evidence/exp3/exp3_recv-sizes-one-message.png}
{Experiment 3 --- four \texttt{recv()} calls returning 6, 10, 7 and 1 bytes, and
the single parsed \texttt{LOGIN} line emitted only after the fourth.}

\shot{/home/vzux/Desktop/A2_evidence/exp3/exp3_four-segments-one-message.png}
{Experiment 3 --- the wire view: four \texttt{[P.]} segments with contiguous
sequence numbers and readable ASCII payloads, then the server's \texttt{OK}.}

\newpage
% ---------------------------------------------------------------------------
\subsection{Experiment 4 --- One Client Should Not Stall the Others}

\textbf{Question.} When Client~1 remains connected but sends no data, can the
Exchange Server still accept and service Client~2? Identify the server operation
that determines the answer.

\textbf{Answer.} \textbf{Yes.} Client~2 connected, sent
\cmd{LOGIN active\_client\textbackslash n} and received \cmd{OK} in
\textbf{0.001\,s} while Client~1 was still holding an unterminated
message.

The operation that decides this is the server's blocking call: it blocks in
\cmd{kevent()} waiting for \emph{readiness across all registered descriptors},
not in \cmd{recv()} on any one socket. \cmd{procstat -k} proves it --- the only
thread's kernel stack is \cmd{\dots\ \_sleep kqueue\_scan kqueue\_kevent
kern\_kevent\_fp kern\_kevent\_generic sys\_kevent \dots}. The server \emph{is}
asleep, but inside \cmd{sys\_kevent}. There is no \cmd{soreceive}, no
\cmd{recv}, nothing tied to Client~1's socket anywhere in the stack. When
Client~2's data arrives the kernel makes fd~6 ready, \cmd{kevent()} returns it,
and the loop services it without revisiting fd~5.

A blocking server that called \cmd{recv(fd5)} in a fixed order would instead be
parked in \cmd{soreceive} on fd~5 indefinitely, and Client~2 would never be
accepted --- precisely the stall this design avoids.

\textbf{Where the incompleteness lives.} Client~1's Recv-Q is \textbf{0, not
20}. The server already performed \cmd{recv()} on fd~5 and received all 20 bytes
(\cmd{[5] recv() -> 20 byte(s)}); they sit in that connection's application-level
\cmd{LineBuffer} awaiting a newline that never comes, so no command line is ever
emitted for fd~5. The kernel socket is drained and quiescent. The partial
message is therefore an \emph{application-layer} state, not a TCP-layer one:
TCP has delivered everything sent and has nothing outstanding.

\textbf{Approach.} Ran the harness, which makes Client~1 send a message with no
terminator and then fall silent, and Client~2 send a complete message
$\approx$2\,s later. Measured the latency the harness reports for Client~2's
reply, then --- while both connections were still open --- inspected the server
from outside: its kernel stack, the socket table, and its descriptor table.

\textbf{Commands and tools.}
\begin{Verbatim}[fontsize=\small,frame=single]
python3 -u experiment.py 4 2>&1 | tee exp4_harness.txt
SRV=$(pgrep -x exchange_server); procstat -k $SRV | tee exp4_kstack.txt
netstat -an -p tcp | grep -E 'Proto|5000'  | tee exp4_netstat.txt
sockstat -4 | grep -E 'USER|exchange'      | tee exp4_sockstat.txt
\end{Verbatim}
\cmd{procstat -k} prints the in-kernel call stack of every thread; this is what
identifies the blocking operation, which timing alone cannot.

\shot{/home/vzux/Desktop/A2_evidence/exp4/exp4_client2-not-stalled.png}
{Experiment 4 --- \texttt{[5] recv() -> 20 byte(s)} with \emph{no} parsed line
following it (Client~1), versus \texttt{[6] recv() -> 20} $\rightarrow$
\texttt{LOGIN active\_client} $\rightarrow$ \texttt{OK} in 0.001\,s (Client~2).}

\shot{/home/vzux/Desktop/A2_evidence/exp4/exp4_procstat-kqueue-not-recv.png}
{Experiment 4 --- the decisive evidence. \texttt{procstat -k} shows the single
thread asleep in \texttt{sys\_kevent}, not in \texttt{soreceive}. Both
connections \texttt{ESTABLISHED} with Recv-Q~0; \texttt{sockstat} shows fds
3/5/6 on one PID.}

\newpage
% ---------------------------------------------------------------------------
\subsection{Experiment 5 --- Multiple Clients and I/O Multiplexing}
\emph{(optional experiment)}

\textbf{Question.} At a particular point during the experiment, which client
connections are actually ready for the server to service, and what evidence from
the running system allows you to determine this?

\textbf{Answer.} Only the connections carrying inbound data were ready:
Client~1 (fd~5, port 65198), Client~3 (fd~7, port 42987) and Client~5 (fd~9,
port 42493). Clients~2 (fd~6) and~4 (fd~8) were fully \cmd{ESTABLISHED} but had
sent nothing, so kqueue never reported them and the server never touched them.

\textbf{Three independent lines of evidence establish this.}

\emph{1. The server's own trace.} The log contains \cmd{recv()} lines for fd~5,
fd~7 and fd~9 only --- 15 bytes each, each followed by the parsed \cmd{LOGIN}.
There is \textbf{no \cmd{[6]} or \cmd{[8]} line anywhere}. The server did not
call \cmd{recv()} on the idle sockets and get \cmd{EAGAIN}; it never called
\cmd{recv()} on them at all. Readiness, not polling, decides what is serviced.

\emph{2. \cmd{netstat} distinguishes them independently.} On the client side of
each connection the three active clients show \textbf{Recv-Q~=~3} --- the
server's 3-byte \cmd{OK} reply sitting unread --- while the two idle clients
show Recv-Q~=~0 in both directions. Even without the server's log, the socket
table alone identifies which connections saw application traffic.

\emph{3. \cmd{procstat} shows the structural cost.} \cmd{procstat -f} lists one
process holding fd~3 (listening), fd~4 (the kqueue) and fds~5--9 (the five
connections). \cmd{procstat -k} shows exactly \textbf{one thread}, asleep in
\cmd{kqueue\_scan}\,/\,\cmd{sys\_kevent}.

The mechanism: all five sockets are registered with \cmd{EVFILT\_READ} on the
same kqueue, and \cmd{kevent()} returns only descriptors the kernel has marked
ready. The loop's work is proportional to the number of \emph{active}
connections, not \emph{open} ones. Idle connections cost a descriptor and some
kernel socket state, but no CPU --- the property that makes the bonus tractable.

\textbf{Approach.} Ran the harness, which opens five connections and then
has clients 1, 3 and 5 send a \cmd{LOGIN} 0.5\,s apart while 2 and 4 stay
silent. During the 15-second observation window, captured three views of the
live server: the socket table, its file-descriptor table, and its thread kernel
stacks. Cross-referenced the fd numbers in the server's accept log against the
ports in \cmd{netstat} to attribute each socket to a specific client.

\textbf{Commands and tools.}
\begin{Verbatim}[fontsize=\small,frame=single]
python3 -u experiment.py 5 2>&1 | tee exp5_harness.txt
netstat -an -p tcp | grep -E 'Proto|5000' | tee exp5_netstat.txt
procstat -f $(pgrep -x exchange_server)   | tee exp5_fds.txt
procstat -k $(pgrep -x exchange_server)   | tee exp5_kstack.txt
\end{Verbatim}

\shot{/home/vzux/Desktop/A2_evidence/exp5/exp5_only-active-clients-serviced.png}
{Experiment 5 --- five \texttt{accepted fd=} lines mapping fd to port, then
\texttt{recv()}/\texttt{LOGIN} for clients 1, 3 and 5 and \textbf{no}
\texttt{[6]} or \texttt{[8]} lines at all.}

\shot{/home/vzux/Desktop/A2_evidence/exp5/exp5_five-conns-one-thread.png}
{Experiment 5 --- \texttt{netstat} (Recv-Q~3 on the three active clients versus
0 on the idle two), \texttt{procstat -f} (fd~3 listen, fd~4 kqueue, fds~5--9
connections on one PID), and \texttt{procstat -k} (one thread in
\texttt{sys\_kevent}).}
\label{fig:e5b}

\newpage
% ---------------------------------------------------------------------------
\subsection{Experiment 6 --- FIN vs.\ RST}

\textbf{Question.} How does an abrupt client termination differ from the orderly
shutdown? Identify the TCP event observed on the network and the corresponding
behaviour of the Exchange Server's socket.

\textbf{Answer.} The two differ in both what crosses the wire and what the
server's socket reports, and the difference is \emph{negotiated} versus
\emph{unilateral}.

\emph{Orderly (Part A, client 127.0.0.1:33130) --- 7 packets.} The client calls
\cmd{shutdown()} on its sending direction, putting a FIN on the wire. The close
is a negotiated four-packet exchange: client FIN \cmd{[F.]} $\rightarrow$ server
ACK \cmd{[.]} $\rightarrow$ server FIN \cmd{[F.]} $\rightarrow$ client ACK
\cmd{[.]}. Each FIN consumes a sequence number and is acknowledged. On the
server the FIN surfaces as \textbf{\cmd{recv()} returning 0} --- a normal
end-of-stream, not an error --- logged as
\cmd{[5] recv() -> 0 : peer closed (FIN)}. The server then closes its own side,
emitting its FIN 142\,$\mu$s later.

\emph{Abortive (Part B, client 127.0.0.1:25854) --- 4 packets.} The client sets
\cmd{SO\_LINGER} with a zero timeout and closes, which makes the kernel send RST
instead of FIN. This is unilateral: \textbf{a single \cmd{[R.]} packet, never
acknowledged}, with no reply from the server and no further exchange of any
kind. On the server the reset surfaces as \textbf{\cmd{recv()} failing with
\cmd{ECONNRESET}}, logged as
\cmd{[5] recv() -> ECONNRESET : peer reset (RST)}. Any queued outbound data
would be discarded rather than delivered, and a subsequent \cmd{send()} would
fail with \cmd{EPIPE}.

So the distinguishing TCP event is the flag itself --- FIN versus RST --- and the
corresponding socket behaviour is \cmd{recv()}~$=$~0 (orderly EOF) versus
\cmd{recv()}~$=$~$-1$ with \cmd{ECONNRESET} (error). The packet counts, 7
versus 4, quantify it: RST needs no acknowledgement and no second direction,
which is exactly why it is ``abrupt''.

\textbf{RST arises from two distinct causes in this capture.} At 16:17:20.509266
the server side sends \cmd{[R.]} in reply to a SYN on port 5000 \emph{before the
server had bound}: a SYN to a port with no listener is answered with RST
(connection refused), never with FIN. That is a different cause from Part~B's
deliberate abortive close at .678706, yet both produce the same single
unacknowledged reset packet. A third RST at .610374 is the harness's readiness
probe closing abortively after confirming the server was up. Only .678706 is the
experiment's own Part~B reset.

\textbf{Note on \cmd{netstat} versus the capture.} The 1-second sampler never
caught \cmd{ESTABLISHED} for Part~A at all --- the client shut down
274\,$\mu$s after connecting, far below the sampling interval. What
the sampler \emph{did} show is the client socket sitting in state \cmd{CLOSED}
from 16:17:21 to 16:17:29, disappearing at 16:17:30. The TCP connection was
fully torn down at .611, but the harness had not yet called \cmd{close()} on the
file descriptor, so the socket still existed --- separating TCP \emph{connection}
state from \emph{socket/fd} lifetime. (This is not itself a FIN-versus-RST
difference: Part~B's fd was closed immediately, so the two are not directly
comparable on that point.)

\textbf{Approach.} Ran the harness, which performs an orderly
\cmd{shutdown()} first and an \cmd{SO\_LINGER\{1,0\}} abortive close second,
with a packet capture and a 1-second state sampler running throughout. Compared
three views of the same two events: the flags on the wire, the server's own
logging of what its \cmd{recv()} call returned, and the socket table. Finally
filtered the capture down to only FIN and RST lines so the two terminations sit
side by side.

\textbf{Commands and tools.}
\begin{Verbatim}[fontsize=\small,frame=single]
su root -c 'tcpdump -l -i lo0 -n -tttt -S "tcp port 5000"' | tee exp6_tcpdump.txt
while true; do date +%T; netstat -an -p tcp | grep 5000; echo; sleep 1; done \
    | tee exp6_states.txt
python3 -u experiment.py 6 2>&1 | tee exp6_harness.txt
grep -E 'Flags \[(F|R)' exp6_tcpdump.txt | tee exp6_fin_rst_only.txt
\end{Verbatim}
\cmd{-S} prints absolute sequence numbers so the FIN's sequence consumption is
visible. The server distinguishes the two paths explicitly at
\cmd{exchange\_server.cpp:382} (\cmd{recv()==0}) and \cmd{:393}
(\cmd{ECONNRESET}).

\shot{/home/vzux/Desktop/A2_evidence/exp6/exp6_server-fin-vs-econnreset.png}
{Experiment 6 --- the application-visible difference: Part~A
\texttt{recv() -> 0 : peer closed (FIN)} versus Part~B
\texttt{recv() -> ECONNRESET : peer reset (RST)}.}

\shot{/home/vzux/Desktop/A2_evidence/exp6/exp6_tcpdump-full-capture.png}
{Experiment 6 --- the complete 17-packet capture covering both terminations.}

\shot{/home/vzux/Desktop/A2_evidence/exp6/exp6_fin-vs-rst-isolated.png}
{Experiment 6 --- only the FIN and RST lines: two FINs (Part~A) and three RSTs
(connection-refused, harness probe, and Part~B's abortive close).}

\shot{/home/vzux/Desktop/A2_evidence/exp6/exp6_states-closed-then-gone.png}
{Experiment 6 --- the client socket in state \texttt{CLOSED} from 16:17:21 to
16:17:29, gone at :30 when the fd is finally closed.}

\newpage
% ---------------------------------------------------------------------------
\subsection{Experiment 7 --- Backpressure and the Slow Receiver}

\textbf{Question.} How does the Exchange Server's TCP connection to the slow
client behave as the client stops reading, and what evidence shows whether this
eventually affects the server's ability to communicate with other clients?

\textbf{Answer.} The connection to the slow client stays \cmd{ESTABLISHED}
throughout and never errors, but undelivered market data accumulates in its
kernel receive buffer, climbing monotonically from 0 to a hard ceiling of
\textbf{85{,}000 bytes}. Over the same period the normal Market-Data client's
Recv-Q never rose above 17 bytes and both Trader connections kept exchanging
orders continuously, with all 5{,}000 trades completing. The slow reader
therefore did \textbf{not} affect the server's ability to serve anyone else ---
the backlog is confined to its own socket.

The mechanism is TCP flow control plus per-connection output queuing. As the
slow client stops calling \cmd{recv()}, its receive buffer fills and the window
it advertises shrinks. The server never blocks on this, because it holds pending
output in a per-connection application-level queue and only calls \cmd{send()}
when kqueue reports that socket writable via \cmd{EVFILT\_WRITE}. Sockets that
are not writable are simply not serviced that iteration; every other descriptor
continues to be.

\textbf{Why Send-Q stays near zero --- the key subtlety.} \cmd{netstat}'s Send-Q
shows only the \emph{kernel} socket buffer. Peak server-side Send-Q toward the
slow client was 34 bytes (Pass~A) and 51 bytes (Pass~B) --- essentially nothing.
That is not evidence that backpressure is absent; it is evidence that the server
never over-commits to a socket the kernel will not accept. The backlog lives in
the server's own userspace output queue, which \cmd{netstat} cannot see. A naive
server calling a blocking \cmd{send()} in a loop would instead fill the kernel
buffer, then block inside \cmd{send()} on the slow client's socket, and every
other client would stall behind it. The near-zero Send-Q is the design working,
not failing.

\textbf{Two passes: a hypothesis tested and rejected.} Pass~A ran with FreeBSD
defaults. Because Recv-Q reached 85{,}000 --- above \cmd{net.inet.tcp.recvspace}
of 65{,}536 --- the initial hypothesis was that receive-buffer auto-tuning
(\cmd{net.inet.tcp.recvbuf\_auto = 1}) was growing the buffer and absorbing the
backlog. Pass~B disabled auto-tuning and re-ran identically. The result was
\textbf{unchanged}: Recv-Q again climbed to exactly 85{,}000 and plateaued, and
Send-Q again stayed negligible. Auto-tuning was therefore \emph{not} the
explanation; the 85{,}000-byte ceiling is a stable property of the receive
socket buffer on this system, and the flat Send-Q is explained by the server's
own write discipline.

\textbf{Approach.} Ran the harness twice with a 1-second \cmd{netstat}
sampler capturing Recv-Q and Send-Q on every connection. Recorded the kernel
buffer limits beforehand so the observed numbers could be compared against the
configured ceilings, and identified which port belonged to which role from the
harness's own connection banner so the sampler output could be interpreted.
Pass~A used the defaults; Pass~B changed exactly one variable
(\cmd{recvbuf\_auto}) to test whether auto-tuning explained the result, and
restored it afterwards.

\textbf{Commands and tools.}
\begin{Verbatim}[fontsize=\small,frame=single]
sysctl net.inet.tcp.sendspace net.inet.tcp.recvspace kern.ipc.maxsockbuf
while true; do date +%T; netstat -an -p tcp | grep -E 'Proto|5000'; echo; \
    sleep 1; done | tee exp7_sendq.txt
python3 -u experiment.py 7 2>&1 | tee exp7_harness.txt
grep -E "Market-Data Client:|Trader:" exp7_harness.txt   # role -> port mapping
su root -c 'sysctl net.inet.tcp.recvbuf_auto=0'          # Pass B only
\end{Verbatim}

\textbf{Measured values.} Buffers: \cmd{sendspace}~32768,
\cmd{recvspace}~65536, \cmd{maxsockbuf}~2097152.
Pass~A (slow~=~40600, normal~=~63111): slow Recv-Q climbed
$102 \rightarrow 544 \rightarrow \dots \rightarrow 85{,}000$; normal 0--17
throughout; server Send-Q to slow peaked at 34.
Pass~B (slow~=~34500, normal~=~44850): slow Recv-Q
$77{,}180 \rightarrow \dots \rightarrow 85{,}000$ then flat; server Send-Q
peaked at 51. Both passes completed all 5{,}000 trades.

\shot{/home/vzux/Desktop/A2_evidence/exp7/exp7_bufsizes.png}
{Experiment 7 --- the configured socket buffer limits.}

\shot{/home/vzux/Desktop/A2_evidence/exp7/exp7_A-slow-recvq-vs-normal.png}
{Experiment 7, Pass A --- three consecutive samples: the slow client's Recv-Q
climbing $84{,}065 \rightarrow 84{,}490 \rightarrow 84{,}932$ while the normal
client sits at 0/17 and both traders keep flowing.}

\shot{/home/vzux/Desktop/A2_evidence/exp7/exp7_A-trades-complete-and-roles.png}
{Experiment 7, Pass A --- live trade traffic, \texttt{Generated 5000 matching
trades.}, the clean four-FIN shutdown, and the role-to-port mapping.}

\shot{/home/vzux/Desktop/A2_evidence/exp7/exp7_B-recvbuf-auto-disabled.png}
{Experiment 7, Pass B --- \texttt{net.inet.tcp.recvbuf\_auto: 1 -> 0}.}

\shot{/home/vzux/Desktop/A2_evidence/exp7/exp7_B-recvq-plateau-85000.png}
{Experiment 7, Pass B --- with auto-tuning disabled the Recv-Q still reaches and
holds exactly 85{,}000 across two consecutive samples, disproving the
auto-tuning hypothesis.}

\shot{/home/vzux/Desktop/A2_evidence/exp7/exp7_B-trades-complete-and-roles.png}
{Experiment 7, Pass B --- \texttt{Generated 5000 matching trades.}
confirming the run completed with auto-tuning disabled, plus the clean
shutdown and role mapping.}

\shot{/home/vzux/Desktop/A2_evidence/exp7/exp7_B-roles.png}
{Experiment 7, Pass B --- role-to-port mapping.}

\newpage
% ---------------------------------------------------------------------------
\subsection{Experiment 8 --- Unexpected Client Disconnection}

\textbf{Question.} When a client disappears unexpectedly while the Exchange
Server is communicating with it, what happens to their TCP connection, and what
network-level evidence allows you to determine how the server detects the
failure?

\textbf{Answer.} The connection is closed \textbf{gracefully}, and the server
detects the failure exactly as it would detect a polite disconnect:
\cmd{recv()} returns~0.

The disappearing client was a \emph{separate process} (server fd~8,
\cmd{127.0.0.1:41329}) killed with \cmd{SIGKILL}, so it had no opportunity to
run shutdown code, close its socket, or send anything deliberately. Nevertheless
a normal \textbf{FIN} appeared on the wire at 17:01:09.148094, and the teardown
was the full four-way exchange: client FIN \cmd{[F.]} $\rightarrow$ server ACK
(.148173) $\rightarrow$ server FIN (.148281) $\rightarrow$ client ACK (.148290).
The server logged \cmd{[8] recv() -> 0 : peer closed (FIN)} and closed fd~8.

The reason is that \textbf{socket teardown is the kernel's responsibility, not
the application's}. When a process dies --- however violently --- the kernel
reclaims its file descriptors, and closing a TCP socket emits FIN. So an abrupt
\emph{process} failure still produces an orderly \emph{connection} close, and at
the socket layer it is indistinguishable from the client calling \cmd{close()}
itself. This is the key contrast with Experiment~6, where RST was produced by
\cmd{SO\_LINGER\{1,0\}} --- an explicit application-level choice about \emph{how}
to close. A crash makes no such choice, so it produces no RST.

Notably, \textbf{306 bytes} of undelivered market data were still sitting unread
in that client's receive buffer at 17:01:08 (the helper process was parked in
\cmd{sleep()}, not reading), and FreeBSD still sent FIN rather than a reset.

RST would arise only afterwards, if the server sent further data to the dead
socket --- the peer's kernel, no longer holding that four-tuple, would answer
with RST. That did not happen here because the event loop saw
\cmd{recv()}~$=$~0 immediately and closed its own side within
187\,$\mu$s, before further market data could be written.

\textbf{The other clients were completely unaffected.} In the same 1-second
sample where the dead connection vanished (17:01:09), the surviving Market-Data
client (57334) and both Traders (53315, 29852) remained \cmd{ESTABLISHED}, and
the capture shows trade traffic to 29852 and 53315 flowing at 17:01:09.077
before the FIN and resuming at .149233 immediately after. The server states it
plainly: \cmd{[-] fd=8 closed (3 client(s) still connected)}.

\textbf{Approach.} Ran the harness, which starts a genuinely separate client
process, generates real trades so the server is actively communicating with it,
then \cmd{SIGKILL}s that process while a second Market-Data client and two
Traders stay connected. Observed the packets around the kill, the socket table
before and after, and the server's own report of what its \cmd{recv()} returned.
Because the killed client is a separate process the harness never prints its
port, so the port had to be recovered from the server's accept log.

\textbf{Commands and tools.}
\begin{Verbatim}[fontsize=\small,frame=single]
su root -c 'tcpdump -l -i lo0 -n -tttt "tcp port 5000"' | tee exp8_tcpdump.txt
while true; do date +%T; netstat -an -p tcp | grep 5000; echo; sleep 1; done \
    | tee exp8_states.txt
python3 -u experiment.py 8 2>&1 | tee exp8_harness.txt
grep -E "Market-Data Client|Trader:|accepted fd=" exp8_harness.txt
\end{Verbatim}

\begin{center}\small
\begin{tabular}{llll}
\toprule
fd & Port & Role & Fate \\
\midrule
5 & 57334 & Surviving Market-Data & survives \\
6 & 53315 & Buyer Trader          & survives \\
7 & 29852 & Seller Trader         & survives \\
8 & 41329 & Disappearing Market-Data & \textbf{SIGKILLed} \\
\bottomrule
\end{tabular}
\end{center}

\shot{/home/vzux/Desktop/A2_evidence/exp8/exp8_roles-and-detection.png}
{Experiment 8 --- fd-to-port mapping including
\texttt{accepted fd=8 from 127.0.0.1:41329}, then the kill moment with
\texttt{[8] recv() -> 0 : peer closed (FIN)} and
\texttt{3 client(s) still connected}.}

\shot{/home/vzux/Desktop/A2_evidence/exp8/exp8_dead-gone-survivors-established.png}
{Experiment 8 --- 41329 present with Recv-Q $187 \rightarrow 306$ at 17:01:07--08,
gone at 17:01:09, while the other three remain \texttt{ESTABLISHED}.}

\shot{/home/vzux/Desktop/A2_evidence/exp8/exp8_tcpdump-kernel-fin.png}
{Experiment 8 --- the complete four-way close of 41329, sandwiched between trade
traffic to 29852/53315 before and after: one frame proving the other clients
kept flowing throughout.}

\newpage
% ===========================================================================
\section{Bonus --- Connection Scalability and I/O Design}
% ===========================================================================

\subsection{Client-generation program}

\cmd{bonus/conn\_generator.cpp} opens and holds $N$ idle TCP connections,
exchanging no application data. It raises its own \cmd{RLIMIT\_NOFILE} to the
hard limit, optionally \cmd{bind()}s each client socket round-robin across a
list of source addresses (see \S\ref{sec:tuple}), reports progress every 5{,}000
connections, and on failure prints the exact \cmd{errno} together with an
interpretation. Helper scripts \cmd{bonus/bonus\_run.sh} (runs one level
end-to-end) and \cmd{bonus/bonus\_measure.sh} (collects every table column in
one pass) accompany it.

\subsection{Resource-measurement table}

All measurements were taken while the stated number of connections were
simultaneously \cmd{ESTABLISHED} and idle.

\begin{center}\small
\begin{tabular}{rrrrrrr}
\toprule
Idle conns & Server RSS & CPU & Server FDs & System-wide FDs & Sock-buf & Max established \\
\midrule
10{,}000 & 5{,}468 KB  & 0.0\% & 10{,}009 & 20{,}102 / 400{,}000  & 0 / 251{,}575 & 10{,}000 \\
20{,}000 & 7{,}432 KB  & 0.0\% & 20{,}008 & 40{,}115 / 400{,}000  & 0 / 251{,}575 & 20{,}000 \\
30{,}000 & 9{,}524 KB  & 0.0\% & 30{,}008 & 60{,}105 / 400{,}000  & 0 / 251{,}575 & 30{,}000 \\
40{,}000 & 11{,}396 KB & 0.0\% & 40{,}009 & 80{,}102 / 400{,}000  & 0 / 251{,}575 & 40{,}000 \\
50{,}000 & 13{,}300 KB & 0.0\% & 50{,}008 & 100{,}105 / 400{,}000 & 0 / 251{,}575 & 50{,}000 \\
60{,}000 & 15{,}616 KB & 0.1\% & 60{,}008 & 120{,}104 / 400{,}000 & 0 / 251{,}575 & 60{,}000 \\
70{,}000 & 17{,}532 KB & 0.1\% & 70{,}009 & 140{,}102 / 400{,}000 & 0 / 251{,}575 & \textbf{70{,}000} \\
\bottomrule
\end{tabular}
\end{center}

\emph{Reading the table.} Server FDs $=$ connections $+\,3$ (listening socket,
kqueue, stdio). System-wide FDs $=$ exactly $2\times$ connections $+\,{\approx}105$:
both endpoints are local, so every connection consumes \textbf{two} file table
entries and \textbf{two} socket structures. This factor of two governs every
system-wide limit in the experiment. Sock-buf is mbuf \emph{clusters} in use
from \cmd{netstat -m}; it is zero at every level (see Q3). The 10k/40k/70k rows
are the values in the submitted screenshots; the others come from separate runs
of the same script, where run-to-run variation is a few KB of RSS.

\subsection{Q1 --- Able to maintain all requested connections?}

Yes at every level from 10{,}000 to 70{,}000 inclusive, \emph{after} tuning; all
seven runs reported \cmd{SUCCESS: all N connections established}.

Before tuning it failed at \textbf{64{,}909 of 70{,}000} with errno~55
(\cmd{ENOBUFS}, ``No buffer space available''). That is the first-failure point
on a default-configured guest; diagnosing it is Q2.

\subsection{Q2 --- First significant bottleneck, with supporting data}
\label{sec:tuple}

\textbf{The kernel's \cmd{socket} UMA zone} --- not memory, CPU, file
descriptors or mbufs. From \cmd{vmstat -z} at the moment of failure:

\begin{Verbatim}[fontsize=\small,frame=single]
ITEM        SIZE   LIMIT     USED    FREE      REQ  FAIL
socket:     1024, 129967,   62006,  67962,  214502,  110
\end{Verbatim}

\begin{itemize}\itemsep2pt
\item \cmd{LIMIT} was 129{,}967 --- identical to the guest's \emph{boot-time}
      \cmd{kern.maxfiles}.
\item \cmd{FAIL 110} is the kernel's own count of refused allocations from that
      zone; the generator observed those same refusals as \cmd{ENOBUFS}.
\item Because both endpoints are local, $N$ connections need $2N$ socket
      structures, so the zone caps connections at $129{,}967/2 = 64{,}983$. The
      run stopped at \textbf{64{,}909} --- within 0.1\% of the prediction.
\end{itemize}

Everything else had large headroom at that moment, which is what rules the
alternatives out: server RSS 16.2\,MB of 4\,GB (0.4\%); 130{,}006 open files of
\cmd{kern.maxfiles} 400{,}000 (32.5\%); 64{,}912 fds of
\cmd{kern.maxfilesperproc} 200{,}000; \textbf{0} mbuf clusters of 251{,}575; CPU
0.0\% with 0.34\,s of CPU time accumulated in total.

\textbf{The non-obvious part.} Raising \cmd{kern.maxfiles} from 129{,}967 to
400{,}000 at runtime did \emph{not} help, because the socket zone is sized once
at boot from that value. The tunable that actually gates socket allocation is
\cmd{kern.ipc.maxsockets}, which \emph{is} writable at runtime and does resize
the zone:
\begin{Verbatim}[fontsize=\small,frame=single]
sysctl kern.ipc.maxsockets=400000     # socket zone LIMIT 129967 -> 400000
\end{Verbatim}

\textbf{A second limit was anticipated and designed around rather than hit.} A
TCP connection is identified by (src IP, src port, dst IP, dst port). Every
connection here targets the same \cmd{127.0.0.1:5000}, so three elements are
fixed and uniqueness rests on the source port alone --- a 16-bit value.
\cmd{net.inet.ip.portrange} defaults to 10000--65535 $=$ 55{,}536 ports, and even
widened to 1024--65535 gives only 64{,}512. \textbf{70{,}000 is therefore
arithmetically unreachable from a single source address.} The generator binds
each client socket round-robin across three loopback aliases (127.0.0.1,
127.0.0.2, 127.0.0.3, added with \cmd{ifconfig lo0 alias}), giving
$3 \times 64{,}512 = 193{,}536$ of tuple space. The server's accept log confirms
all three addresses in use throughout the run.

\textbf{A third, softer symptom} appeared at ${\approx}65{,}000$ connections: the
guest stopped accepting new inbound SSH sessions entirely, and recovered the
instant the connections were released. Existing sessions were unaffected.

\subsection{Q3 --- How the concurrency/I/O design contributes}

The server is a single-threaded \cmd{kqueue()} event loop; it creates no thread
and no process per connection, and \cmd{procstat -k} confirms one thread
throughout (Experiments 4 and 5).

The consequence is visible directly in the table. Server RSS rises from
5{,}468\,KB at 10{,}000 connections to 17{,}532\,KB at 70{,}000 --- a marginal
cost of $\approx$\,2{,}012\,KB per 10{,}000 connections, i.e.\ \textbf{about 201
bytes per idle connection}, on a 3.4\,MB baseline, and the relationship is linear
across all seven levels. CPU stays at 0.0--0.1\% and accumulates 0.33\,s while
holding 70{,}000 connections, because idle sockets are never ready and so are
never visited by the loop.

The design therefore contributes to the bottleneck only indirectly: it pushes
per-connection application cost so low that the application is nowhere near
being the limit, and what fails first is a \emph{kernel} allocator ceiling. That
is the intended outcome --- the wall is an OS resource, not the program.

\textbf{Supporting finding.} mbuf cluster usage is \emph{identical}
(320/696/1016 4k jumbo clusters) at 1{,}000 connections and at 70{,}005. Idle
sockets allocate \textbf{no socket buffer memory at all}; \cmd{sendspace} and
\cmd{recvspace} are ceilings applied as data arrives, not reservations made at
connect time. This is the same property that kept Send-Q near zero in
Experiment~7, observed from the opposite direction.

\subsection{Q4 --- Would changing the I/O mechanism help?}

No --- every alternative is worse, and the measurements say why.

\textbf{Thread-per-connection:} 70{,}000 threads, each with its own stack. Even
at a conservative 512\,KB per stack that is $\approx$34\,GB of address space
against 4\,GB of RAM; at FreeBSD's default it is far more. The measured
per-connection cost here is $\approx$201 bytes --- roughly four orders of
magnitude smaller. Thread-per-connection could not reach 70{,}000 on this
machine at any tuning, and would add context-switch and scheduler cost that
kqueue avoids entirely.

\textbf{\cmd{select()}:} bounded by \cmd{FD\_SETSIZE} (1024 by default) and
requires rebuilding and scanning the descriptor set on every call. Unusable at
this scale.

\textbf{\cmd{poll()}:} no \cmd{FD\_SETSIZE} limit, but the caller passes the
entire array on every call and the kernel scans all of it --- $O(n)$ per
iteration in the number of \emph{open} descriptors. At 70{,}000 mostly-idle
connections that is 70{,}000 entries examined to find the handful that are
ready. \cmd{kqueue} is $O(\mathrm{ready})$: registration happens once, and
\cmd{kevent()} returns only what changed.

Since the observed limit is a kernel socket-zone ceiling rather than anything
the application controls, no change of I/O mechanism would raise it. Only the
sysctl does; changing the mechanism could only make the application side worse.

\subsection{Q5 --- Quantified trade-off after a change}

One configuration variable changed, everything else held constant:

\begin{center}\small
\begin{tabular}{rrrl}
\toprule
\cmd{kern.ipc.maxsockets} & socket zone LIMIT & Max connections & Result \\
\midrule
129{,}967 (boot value) & 129{,}967 & 64{,}909 & \cmd{ENOBUFS} at 64{,}909 \\
400{,}000              & 400{,}000 & \textbf{70{,}005} & SUCCESS, headroom remains \\
\bottomrule
\end{tabular}
\end{center}

Raising the zone from 129{,}967 to 400{,}000 lifted the achievable connection
count from 64{,}909 to at least 70{,}005 --- an 8\% increase, for a change that
costs no memory until the sockets are actually created. At 70{,}005 connections
the zone held 140{,}104 of 400{,}000 entries (35\%), so the new ceiling is
roughly 200{,}000 connections: the limit has moved well beyond the assignment
target.

The trade-off is that the zone is a pre-declared kernel ceiling --- raising it
permits more kernel memory to be committed to sockets if they are created. The
measured cost of actually reaching 70{,}000 was 17.1\,MB of server RSS and
140{,}102 file table entries, both trivial against 4\,GB of RAM, which is why
the increase is essentially free on this system.

\textbf{A second change was considered and rejected on evidence:} reducing
\cmd{net.inet.tcp.sendspace}/\cmd{recvspace} to shrink per-connection buffer
memory. The measurements show idle sockets allocate zero mbuf clusters, so that
change would have measured no difference. It is recorded here because ruling it
out is what established that socket-buffer memory is not on the critical path.

\subsection{Reproducing the tuning}

Runtime-only; re-apply after any reboot (as root):
\begin{Verbatim}[fontsize=\small,frame=single]
sysctl kern.ipc.maxsockets=400000    # the important one -- sizes the socket zone
sysctl kern.maxfiles=400000
sysctl kern.maxfilesperproc=200000
sysctl kern.ipc.somaxconn=4096
sysctl net.inet.ip.portrange.first=1024
ifconfig lo0 alias 127.0.0.2/32
ifconfig lo0 alias 127.0.0.3/32
\end{Verbatim}
Then: \cmd{sh bonus/bonus\_run.sh <count>} for one level, or
\cmd{sh bonus/bonus\_measure.sh 5000} to measure a running server.

\subsection{Bonus screenshots}

\shot{/home/vzux/Desktop/A2_evidence/bonus/bonus_10000_measurements.png}
{Bonus --- 10{,}000 idle connections: all measurements plus
\texttt{SUCCESS: all 10000 connections established}.}

\shot{/home/vzux/Desktop/A2_evidence/bonus/bonus_40000_measurements.png}
{Bonus --- 40{,}000 idle connections.}

\shot{/home/vzux/Desktop/A2_evidence/bonus/bonus_70000_measurements.png}
{Bonus --- 70{,}000 idle connections: 17.1\,MB RSS, 0.1\% CPU, 70{,}009 server
fds, 140{,}102 of 400{,}000 system-wide files, zero mbuf clusters.}

\shot{/home/vzux/Desktop/A2_evidence/bonus/bonus_socket-zone-limit-evidence.png}
{Bonus --- the bottleneck evidence. The \texttt{socket} UMA zone now at
\texttt{LIMIT 400000}, with \texttt{FAIL 110} still recorded from when the limit
was 129{,}967.}

\end{document}
